%% LyX 2.4.4 created this file.  For more info, see https://www.lyx.org/.
%% Do not edit unless you really know what you are doing.
\documentclass[oneside,english]{book}
\usepackage[T1]{fontenc}
\usepackage[latin9]{inputenc}
\setcounter{secnumdepth}{3}
\setcounter{tocdepth}{3}
\usepackage{amsmath}
\usepackage{amsthm}
\usepackage{amssymb}
\usepackage{geometry}
\geometry{verbose,tmargin=1in,bmargin=1in,lmargin=1.5in,rmargin=1.5in}

\makeatletter
%%%%%%%%%%%%%%%%%%%%%%%%%%%%%% Textclass specific LaTeX commands.
\theoremstyle{definition}
\newtheorem{defn}{\protect\definitionname}
\theoremstyle{plain}
\newtheorem{prop}{\protect\propositionname}
\ifx\proof\undefined\
  \newenvironment{proof}[1][\proofname]{\par
    \normalfont\topsep6\p@\@plus6\p@\relax
    \trivlist
    \itemindent\parindent
    \item[\hskip\labelsep
          \scshape
      #1]\ignorespaces
  }{%
    \endtrivlist\@endpefalse
  }
  \providecommand{\proofname}{Proof}
\fi
\theoremstyle{plain}
\newtheorem{thm}{\protect\theoremname}

%%%%%%%%%%%%%%%%%%%%%%%%%%%%%% User specified LaTeX commands.
\usepackage{hyperref}
\usepackage[usenames,dvipsnames]{xcolor} %adds additional color options
\hypersetup{colorlinks=true, urlcolor=Blue}%Blue

\usepackage{multicol}

\usepackage{tikz}
\usepackage{tikz-3dplot}
\usetikzlibrary{calc,arrows,shapes,backgrounds}

\usetikzlibrary{patterns}
\usetikzlibrary{plotmarks}
\usepackage{pgfplots}
\pgfplotsset{compat=newest}

\usepackage{calc}
\usepackage[nomessages]{fp}

\usepackage{datetime}
%\usepackage{subcaption}

%\usepackage{amsthm}
% Added by lyx2lyx
\setlength{\parskip}{\smallskipamount}
\setlength{\parindent}{0pt}

\makeatother

\usepackage{babel}
\providecommand{\definitionname}{Definition}
\providecommand{\propositionname}{Proposition}
\providecommand{\theoremname}{Theorem}

\begin{document}

\chapter{\label{chap:Rings}Rings}

\section{Rings and Subrings}

We now continue to generalize the familiar arithmetic operations of
addition and multiplication. Just as with groups, results from the
algebra of arithmetic will be extended to a host of non-numerical
objects, e.g., matrices, functions, polynomials, etc.
\begin{defn}
A \emph{ring} is a nonempty set, say $R$, with two binary operations
(usually written as multiplication and addition)\footnote{\emph{But of course the operations are not necessarily addition and
multiplication. $+$ and $\times$; $\oplus$ and $\otimes$; are
among the different notations that may be used in varying texts. }} that satisfies the following:

\begin{enumerate}
\item If $a,b\in R$ then $a+b\in R$. \hfill(closed under addition)
\item $a+(b+c)=(a+b)+c$ \hfill(associative addition)
\item $a+b=b+a$ \hfill(commutative addition)
\item There exists an $0_{R}$ such that $a+0_{R}=0_{R}+a=a$ \hfill(additive
identity)
\item There exists an $-a\in R$ such that $a+(-a)=0_{R}$\\
for every $a\in R$. \hfill(additive inverse)
\item If $a,b\in R$ then $ab\in R$. \hfill(closed under multiplication)
\item $a(bc)=(ab)c$ \hfill(associative multiplication)
\item $a(b+c)=ab+ac$ and $(a+b)c=ac+bc$ \hfill(distribution holds).
\end{enumerate}
\end{defn}
\begin{itemize}
\item Note that parts 1-5 say that a ring is an \emph{abelian group.}\footnote{Commutative groups are also called abelian in honor of Norwegian mathematician
Niels Abel.}\emph{ }
\end{itemize}
\begin{prop}
If $a$ and \textbf{$b$} are two elements in $R$ then $(-a)(-b)=ab$.
\end{prop}
\begin{proof}
By definition $a+(-a)=0$. $(a+(-a))(-b)=0$ implies $a(-b)+(-a)(-b)=0$
by distribution. The additive inverses of $a(-b)$ is unique thus
$a(-b)+(-a)(-b)=0$ implies that $(-a)(-b)=ab$.
\end{proof}
\begin{itemize}
\item In particular for rings like $\mathbb{Z}$ and $\mathbb{R}$, the
product of two negatives is a positive. Note the importance of distribution.
\end{itemize}
\begin{defn}
A \emph{commutative ring }is a ring $R$ such that $ab=ba$ for all
$a,b\in R$.
\end{defn}

\begin{defn}
A ring with \emph{identity}\footnote{The term \emph{unity} is often used in place identity. A ring with
an identity element is said to ``have unity''.}\emph{ }is a ring $R$ with an element $1_{R}$ such that $a1_{R}=1_{R}a$
for all $a\in R$.
\end{defn}
\begin{itemize}
\item It is important to check that a potential $1_{R}$ is both a \emph{right
identity, $a1_{R}=a$}; and a \emph{left identity}, $1_{R}b=b$.
\end{itemize}
\begin{thm}
If a ring has an identity, it is unique. 
\end{thm}
\begin{proof}
Suppose to the contrary that a ring $R$ has two identity elements
$1$ and $1'$. $a1=a$ for all $a\in R$ including $1'\in R$. Thus
$1'\cdot1=1$' and $1\cdot1'=1$ . By the definition of a ring $1'=1\cdot1'=1'\cdot1=1$.
Therefore $1'=1$, and the multiplicative identity is unique.
\end{proof}

\section{Subrings}
\begin{defn}
Let $R$ be a ring and $S$ be a subset of $R$.\footnote{Notated: $S\subseteq R$.}
If $S$ is a ring under the same operations as $R$, then $S$ is
a \emph{subring}.
\end{defn}
The following two theorems provide shortcuts when trying to show that
a subset is a subring. Particularity, The second theorem can save
quite a bit of work.
\begin{thm}
Suppose that $R$ is a ring and $S\subseteq R$ such that the following
hold:

\begin{enumerate}
\item If $a,b\in S$ then $a+b\in S$ (closed under addition)
\item If $a,b\in S$ then $ab\in S$ (closed under multiplication)
\item $0_{R}\in S$.
\item If $a\in S$ then there is an $-a\in S$ such that $a+(-a)=0_{R}$.
\end{enumerate}
Then $S$ is a subring of $R$.
\end{thm}
\begin{proof}
Parts 2, 3, 7, and 8 of the Ring definition hold for all elements
in $R$ since it is a ring. So they hold for all elements in $S$.
The remaining requirements 1-4 in the theorem satisfy the remaining
requirements. 
\end{proof}
\begin{itemize}
\item Rings and thus subrings must be nonempty. A requirement which is covered
by 3. Since $0_{R}\in S\Rightarrow$ $S\neq\emptyset$. 
\end{itemize}
\begin{thm}
\label{thm: subring test}A nonempty subset $S$ of a ring $R$ is
a subring if whenever $a,b\in S$:

\begin{enumerate}
\item $ab\in S$ (closed under multiplication)
\item $a-b\in S$ (closed under subtraction)
\end{enumerate}
\end{thm}
\begin{proof}
If $a-b\in S$ for all $a,b\in S$ then $S$ is a subgroup of $R$
(subgroup test), and since $R$ is additively commutative it follows
that $S$ is a abelian subgroup of $R$ under addition (this also
means that $0_{R}$). $S$ is multiplicatively and additively associative;
and $S$ is distributive over addition because these properties hold
in $R$ .
\end{proof}

\section{Integral Domains and Fields}

For $\mathbb{Z}_{n}$ it is easy to see that $ab=0$ does \emph{not
}imply that $a$ or $b$ is $0$. For example, in $\mathbb{Z}_{6}$
$2\cdot3\mod 6=6\mod 6=0$. Of course many rings (indeed the most
familiar rings) do have this property. In $\mathbb{R}$, if $x\cdot y=0$
then $x=0$ or $y=0$. 
\begin{defn}
An \emph{integral domain }is a commutative ring $R$ with identity
such that $1_{R}\neq0_{R}$\footnote{The condition $1_{R}\neq0_{R}$ is necessary so that the ring $\{0\}$
is not an integral domain. } and if $ab=0_{R}$ then $a=0_{R}$ or $b=0_{R}$ for any $a,b\in R$.\footnote{If $a,b$ are nonzero elements in a ring and $ab=0$ then $a$ and
$b$ are called \emph{zero divisors}. Thus an integral domain is a
commutative ring with identity that has no zero divisors. }
\end{defn}
\begin{thm}
\label{thm:cancel-ID}Let $R$ be an integral domain and let $a,b,c\in R$
where $c\neq0$. If $ac=bc$, then $a=b$.
\end{thm}
\begin{proof}
Since rings have distribution, $ac=bc\Rightarrow ac-bc=0\Rightarrow(a-b)c=0$.
Because $R$ is an integral domain $a-b=0$ which implies that $a=b$.
\end{proof}
Note that we could not prove theorem \ref{thm:cancel-ID} using division,
i.e., $\frac{ac}{c}=\frac{bc}{c}$, because we cannot assume that
$c^{-1}\in R$. For example, $3x=3y\Rightarrow x=y$ in $\mathbb{Z}$
but $\frac{1}{3}\notin\mathbb{Z}$. While $\mathbb{Z}$ is an integral
domain, it has no element $x$ such that $3x=1$. That is, $3$ (and
all other numbers except $1$) does not have a multiplicative inverse.
All the elements in $\mathbb{Q}$, $\mathbb{R}$, and $\mathbb{C}$
(except $0$) do have multiplicative inverses. Thus $\mathbb{Q}$,
$\mathbb{R}$, and $\mathbb{C}$ are all examples of the following
definition.
\begin{defn}
A \emph{field }is a commutative ring $R$ with identity $1_{R}\neq0_{R}$
\footnote{Once again, this assumption is needed to that $\{0\}$ is excluded
from this class.} if for every $a\neq0_{R}$ in $R$, there exists an $a^{-1}\in R$
such that $aa^{-1}=1_{R}$. 
\end{defn}
\begin{itemize}
\item $\mbox{Sets}\supseteq\mbox{Groups}\supseteq\mbox{Rings}\supseteq\mbox{Commutative Rings with identity}\supseteq\mbox{Intergal Domains}\supseteq\mbox{Fields}$
\end{itemize}
\begin{thm}
$\mathbb{Z}_{p}$ is a field if $p$ is a prime.
\end{thm}
\begin{proof}
This is left as an exercise.  
\end{proof}

\begin{thm}
\textbf{\emph{Cancellation. }}Let $a,b,c$ be elements in an integral
domain. Suppose that $a\neq0$. Then $ab=ac\Rightarrow b=c$.
\end{thm}
\begin{proof}
Let $a\neq0$ for elements $a,b,c$ in an integral domain, $R$. $ab=ac\Rightarrow a(b-c)=0_{R}$.
Because $R$ is an integral domain either $a=0$ or $b-c=0$. By assumption,
$a\neq0$ so then $b-c=0\Rightarrow b=c$.
\end{proof}
\begin{thm}
A field is an integral domain.\footnote{Often a \emph{field }is defined as an integral domain in which every
nonzero element has an multiplicative inverse.}
\end{thm}
\begin{proof}
This is left as an exercise. 
\end{proof}
\begin{thm}
A finite integral domain is a field. 
\end{thm}
\begin{proof}
Let $R$ be a finite integral domain with identity. Let $a\in R$
and $a\neq0_{R}$. All we need to show is that $a$ has an inverse.
If $a=1_{R}$ then $a^{-1}=1_{R}$. Now suppose that $a\neq1_{R}$
(if no such exists then we are done). Since $R$ is finite there are
two integers such that $i<j$ and $a^{i}=a^{j}$ (otherwise $R$ would
be infinite). By cancellation (above theorem) we have that $a^{j-i}=1_{R}$.
Because $a\neq1_{R}$ we know that $j-i>1$. Thus $a^{j-i-1}(a)=1_{R}$
as needed.
\end{proof}

\end{document}
